\section{最大公因数、最小公倍数与 Bézout 恒等式}


\subsection{定义与存在性}


对不全为零的整数 $a,b$，定义

\[
\gcd(a,b)=\max\{d\in\mathbb N_{>0}:d\mid a\text{ 且 }d\mid b\}.
\]

这个集合含 $1$；若 $a\ne0$，其中每个元素不超过 $|a|$，所以最大值存在。若 $a=0$，用 $|b|$ 作上界即可。另约定 $\gcd(0,0)=0$；它不由上述取最大值的式子定义。

对 $ab\ne0$，定义

\[
\operatorname{lcm}(a,b)=\min\{m\in\mathbb N_{>0}:a\mid m\text{ 且 }b\mid m\}.
\]

集合非空，因为它含 $|ab|$。任一参数为 $0$ 时，正公倍数集合为空，本稿另约定 $\operatorname{lcm}(a,0)=0$。

\subsection{最小正线性组合等于 gcd}


设 $(a,b)\ne(0,0)$，令

\[
h=\min\{au+bv:u,v\in\mathbb Z,\ au+bv>0\}.
\]

集合非空：至少有一个非零参数，改变其系数的符号即可得到正数。由自然数良序，$h$ 存在。

\textbf{第一步：$h$ 是公因子。} 取 $h=au_0+bv_0$，对 $a$ 作带余除法：

\[
a=qh+r,\qquad 0\le r<h.
\]

于是 $r=a(1-qu_0)+b(-qv_0)$ 也是整数线性组合。若 $r>0$，就与 $h$ 的最小性矛盾，故 $r=0$，即 $h\mid a$。同理 $h\mid b$。

\textbf{第二步：所有公因子都整除 $h$。} 若 $d\mid a,b$，则 $d\mid au_0+bv_0=h$。特别地，所有正公因子都不超过 $h$。结合第一步，得到

\[
\gcd(a,b)=h=au_0+bv_0.
\]

这就是 \textbf{Bézout 恒等式}。它还说明 gcd 是整除偏序中的最大公因子：任意公因子都整除 gcd，而不只是数值上不大于它。$(a,b)=(0,0)$ 时，等式 $0=0u+0v$ 仍成立，但不存在最小正线性组合。

\subsection{两个常用推论}


\textbf{互素判据。} $\gcd(a,b)=1$ 当且仅当存在整数 $u,v$ 使 $au+bv=1$。

\textbf{互素的乘积性质。} 若 $\gcd(a,b)=\gcd(a,c)=1$，取 $au+bv=1$、$as+ct=1$，相乘得

\[
1=a(aus+uct+bvs)+bc(vt),
\]

故 $\gcd(a,bc)=1$。这条性质后来用于 CRT 的构造。
